\section{概述}
\label{sec:intro}

大型转动设备是电力、船舶、航空等领域的核心动力装备。以火电厂
为例，汽轮机、发电机、给水泵等设备往往由多段转子串联组成一条长轴系，各段转子
之间通过联轴器连接。若相邻两段轴线之间存在偏移或折角，机组在运行时便会产生附加
的交变应力与振动，轻则加速轴承与密封件磨损，重则引发动静碰摩甚至断轴事故。因此，
轴系对中是设备安装与检修中的一道关键工序：需要精确测量待对中轴上若干采样点相对
某条理想基准轴的位置偏差，并据此调整设备的抬升与平移。汽轮机转子的直径通常在
1--2\,m 之间，工程上对相邻轴段的对中精度要求达到两丝，即 $20\,\mu\mathrm{m}$
量级，这对测量方法的分辨力与稳定性都提出了很高的要求。

目前行业内仍以钢丝找中法为主。其做法是在待对中的轴系中间张紧一根钢丝作为直线
基准，再用百分表或电接触方式逐点测量轴表面到钢丝的距离，从而还原各采样点相对
基准的偏差。这一方法原理直观、器材简单，但在现场使用中暴露出若干难以回避的缺陷。
其一，钢丝基准无法长期保留：每完成一轮测量都要拆除，隔日复测或调整之后复核时又
需重新架设与张紧，频繁的拆装既拖慢工期，也会干扰同一现场其他工序的并行施工。
其二，横贯现场的张紧钢丝处于人员通行的空间高度，容易绊拌施工人员，构成安全隐患。
其三，也是对精度影响最本质的一点，钢丝在自重作用下会产生下坠挠度，其垂度沿跨度
呈非线性分布并随张力波动，还对气流与机械振动敏感，使这条“基准直线”本身并不真正
是一条直线，从而在毫米量级引入难以彻底修正的系统偏差。

用一条激光束替代钢丝，是消除下坠挠度、免除反复拆装的自然思路。然而问题的难点并
不在于产生一条激光，而在于如何精确测量“空间中任意目标点到这条激光基准轴的垂直
距离”。判断一个已有方法是否可用，关键只看它能否（直接或间接）给出这一几何量，
而与其内部是否用到激光、激光扮演何种角色无关。以此为准，现有可比方案各有硬伤：
精度极高的激光跟踪仪须逐点安放合作反射器、成本高且效率低~\cite{api_radian,leica_at960}；
现役的位置敏感探测器靶板对中仪只能测激光斑击中探测器的那一单点~\cite{easylaser_xt770,fixturlaser_nxa}；
学术上最接近的双目激光线重建方法虽语义一致，精度却仅停留在毫米级~\cite{zhong2021line3d}。
这些方法与本文场景的具体差异将在第~\ref{sec:related} 节详细讨论。

针对上述空白，本文提出 DCPAM（Dual-Camera-Based Point-to-Axis Measurement）
测量模型。其核心思想是：用一条激光束充当替代钢丝的直线基准，在激光光路的前、后
两处各布置一个相机模块，每个模块通过一片 $45^\circ$ 倾斜的反射镜把激光引向一块
毛玻璃接收屏，相机隔着取景框拍摄激光在毛玻璃上漫散射形成的光斑。由两个模块拍到
的光斑分别反投影、并借助各自取景框的位姿标定统一到设备坐标系，再关于反射镜做
镜面反射得到激光的两个虚像点；这两个虚像点确定了激光基准轴在设备坐标系下的三维
空间直线，进而可对固定在测杆末端的目标点解析地计算其到该直线的垂直距离。整个
测量以非接触方式完成，激光基准不存在下坠问题，也无需在每次测量时重新架设。

本文的主要贡献可归纳为三点。首先，明确提出并形式化了“空间任意点到激光基准轴
垂距”这一在现有商业产品与文献中尚未被直接解决的测量问题，并系统梳理了相关方法
在该问题上的适用边界。其次，建立了基于双相机、接收屏与镜面虚像重建的完整几何
测量模型，给出从像素光斑到点线距离的端到端计算链路与相应的标定方法。最后，围绕
重复精度与误差传播开展了实验与精度分析，定位了当前系统的主导误差来源并指出
改进方向。本文其余部分组织如下：第~\ref{sec:related} 节梳理相关测量方法并定位
本文的空白；第~\ref{sec:model} 节给出 DCPAM 的坐标系约定、测量模型与标定方法；
第~\ref{sec:exp} 节报告实验结果与误差分析；第~\ref{sec:conclusion} 节总结全文
并展望后续工作。
